Actualmente gracias al gran desarrollo de las capacidades de computo y el desarrollo de bases de datos amplias, existen tecnologías que dan muy buenos resultados, en el contexto de los métodos anteriormente mencionados, los métodos que aparecen y reinantes de la actualidad corresponderían a la categoría de modelos guiados por datos.

Algunos de los más interesantes:

\begin{itemize}
    \item \textbf{Open-Unmix: } implementación de red neuronal referente en el campo de la separación musical. Provee de modelos preparados para su uso que permite separacion de pistas pop en cuatro fuentes: vocales, baterías, bajo y resto. Se trata de un modelo abierto basado en deep learning y que permite reproducibilidad \cite{openunmix2019}.
    \item \textbf{Spleeter: } conjunto de modelos pre-entrenados de separación de fuentes en pista musicales. Basados en redes neuronales convolucionales con salto de conexiones ( que se traducen en codificadores-decodificadores llamados \textit{U-Nets}) de doce capas, utilizando función de pérdida \textit{L1-norm}, entrenados con bases de datos internas de \textit{Deezer} y con espectrogramas de la mezcla como entrada \cite{Deezer}. 
    \item \textbf{Demus / Hybrid Demucs: } modelo avanzado que, en función del tipo de pista de entrada, toma como entrada su espectrograma, su onda, o ambos. De los estados más modernos en el campo \cite{HybridDemucs}.
\end{itemize}

Además, mencionar también \textbf{MUSDB18} que es un dataset/benchmark de referencia que contiene 150 conjuntos de mezcla y cuatro fuentes comprendiendo bajo, batería, voz y resto. Permite comparar sistemas y está dividido en entrenamiento y test. Ha sido el \textit{dataset} utilizado en este trabajo de fin de grado.